\section{总结与延伸}

\subsection{核心结论}

\begin{importantbox}{十二条可迁移结论}
\begin{enumerate}
  \item AI 浪潮只有穿过应用扩散和组织重构，才产生生产率。
  \item Useful compute 受 semiconductor、memory、interconnect 与 facility 的共同约束。
  \item Memory wall 包含 capacity、bandwidth 与 communication 三种失配。
  \item SRAM、DRAM、HBM 与 SSD 必须按 working-set 温度分层使用。
  \item PUE 只衡量 facility overhead；utilization 与 workload efficiency 同样关键。
  \item Grid、facility、accelerator 与 model demand 的时钟不同，modularity 用来管理 regret。
  \item Generative、agentic 与 physical AI 需要不同 evidence 和 accountability。
  \item 医疗搜索加速不能跳过 clinical governance，机器人手术也不等于 autonomous surgery。
  \item Model agnostic 是 evaluation-driven portfolio，不是取消标准。
  \item Capacity 必须进入 utilization--cost--adoption flywheel 才是 productive asset。
  \item Data sovereignty 要同时设计 location、control、jurisdiction、portability 和 recovery。
  \item Government Agent 的决定性难点在 workflow、data、human-in-the-loop 与经济责任。
\end{enumerate}
\end{importantbox}

\subsection{实践作业：设计一个国家级 AI capacity-to-outcome 系统}

选择 healthcare、energy、education 或 public administration 之一，完成以下设计：

\begin{enumerate}
  \item 画出 grid power 到 useful outcome 的资源会计链，并定义 PUE、utilization 与 workload efficiency；
  \item 估算 working set、memory traffic、network traffic 和 latency SLO，判断 compute/memory/network 瓶颈；
  \item 定义 training 与 inference 的三种 demand scenario，以及设施模块化和采购策略；
  \item 建立 model-agnostic evaluation 与 routing policy，包含 fallback 和 version rollback；
  \item 设计 data sovereignty 五维控制，并给出 region/provider failure 演练；
  \item 把 agent 从 task 扩展到 workflow，列出 state、permission、exception、appeal 和 human gate；
  \item 定义 adoption、service quality、equity 和 total cost 指标，说明何时停止或缩减项目。
\end{enumerate}

\begin{knowledgebox}{验收标准}
高质量方案不会只写“购买 GPU、部署大模型、建设平台”。它应把不可逆的 power/facility 决策与可替换的 hardware/model 分开，把讲者估计与验证事实分开，并给出从 capacity 到 public outcome 的可测中间变量。
\end{knowledgebox}

\subsection{拓展阅读}

\begin{enumerate}
  \item Gholami et al., \href{https://arxiv.org/abs/2403.14123}{AI and Memory Wall}。
  \item Sebastian et al., \href{https://www.nature.com/articles/s41565-020-0655-z}{Memory devices and applications for in-memory computing}。
  \item Saudi Digital Government Authority, \href{https://dga.gov.sa/sites/default/files/2024-08/Cloud%20Computing%20and%20its%20Role%20in%20Accelerating%20Digital%20Transformation%20and%20its%20Sustainable%20Impact%20in%20the%20Government%20Sector%20-%20V1.0.pdf}{Cloud Computing and Digital Transformation}。
  \item KFSHRC, \href{https://www.kfshrc.edu.sa/en/news/2024/09/kfshrc-pioneers-worlds-first-fully-robotic-heart-transplant}{World's first fully robotic heart transplant}。
  \item Groq, \href{https://groq.com/news_press/groq-and-aramco-digital-announce-partnership-to-establish-worlds-largest-inferenc-data-center-in-saudi-arabia}{Aramco Digital inference partnership} 与 \href{https://groq.com/news_press/saudi-arabia-announces-1-5b-expansion-to-fuel-ai-powered-economy-with-groq}{later LEAP 2025 expansion}。
  \item Saudi Data and AI Authority, \href{https://sdaia.gov.sa/en/SDAIA/about/Pages/Strategy.aspx}{National Strategy for Data \& AI}。
\end{enumerate}

\end{document}
